\scanpage{234}
sss les rayons
sont incidents.

\mspage{119 bis}
\uline{NB} La notion de prégraphe $\Gamma=(S,S^\vee,R)$ avec groupe $Z$ d'homothéties virtuelles
\[
 Z\subset\mathrm{Aut}(\Gamma),\qquad Z\text{ ss-gr. comm.} \tag*{(48.1)}
\]
satisfaisant les conditions
\[
 \left\{\begin{aligned}
 &\text{a) }Z\text{ opère librement sur }S\text{ et sur }S^\vee\\
 &\text{b) }\deleted{\unsure{Z}}\ \text{On a la condition (41) (ou (41')).}
 \end{aligned}\right.\tag*{(48.2)}
\]
peut s'expliciter \unsure{aussi|ainsi} \deleted{\unsure{aisém.}}. On peut identifier, à cause
de a), \deleted{\unsure{S}} $S$ et $S^\vee$ comme des $Z$-torseurs relatifs
sur deux ens $S^{\unsure{\leftrightharpoons|\natural|\leftrightarrows|\flat}}$ et $S^{\vee\leftrightharpoons}$, liés par une relation
\diagram{\begin{tikzcd}[ampersand replacement=\&,column sep=small]
{\text{(48.3)}} \& \& R^\leftrightharpoons\arrow[dl,swap,"{\unsure{p_1|p}}"]\arrow[dr,"p^\vee"] \&
 {\subset\underset{\text{ens des rayons}}{S^\leftrightharpoons}\times\underset{\text{ens des corayons}}{S^{\vee\leftrightharpoons}}} \&
 {\begin{gathered} (R^\leftrightharpoons\text{ la ``relation d'in-}\\\deleted{\text{relation d'incidence}}\\
 \text{cidence'' entre}\\\text{rayons et corayons)}\end{gathered}}\\
\& S^\leftrightharpoons \& \& S^{\vee\leftrightharpoons} \&
\end{tikzcd}}
et on a un isom de $Z$-torseurs
\[
 p^*(S/S^\leftrightharpoons)\simeq p^{\vee *}(S^\vee/S^{\vee\leftrightharpoons})
 \tag*{(48.4)}
\]
-- la relation $R$ entre $S$ et $S^\leftrightharpoons$ s'identifie à
$p^*(S/S^\leftrightharpoons)$ et $p^{\vee *}(S^\vee/S^{\vee\leftrightharpoons})$ identifiés par (44.4), s'envoyant
dans $S$ et dans $S^{\deleted{\unsure{*}}\vee}$ via les diagrammes commutatifs
\diagram{\begin{tikzcd}[ampersand replacement=\&,column sep=small]
\& S\arrow[d] \& {\deleted{\unsure{S^\leftrightharpoons}}}\arrow[d] \&
{\overbrace{p^*(S/S^\leftrightharpoons)\simeq p^{\vee *}(S^\vee/S)}^{R}}
\arrow[ll]\arrow[d,start anchor=south west,end anchor=north west]\arrow[d,start anchor=south east,end anchor=north east]\arrow[rr,<->]
\& {\deleted{\unsure{S}}} \& S^\vee\arrow[d]\\
{\text{(48.5)}} \& S^\leftrightharpoons \& {\deleted{\unsure{R^\leftrightharpoons|R'^{\leftrightharpoons}}}} \&
R^\leftrightharpoons\arrow[ll,"{\unsure{p_1|p_4|p_{\flat}}}"']\arrow[rr,swap,"p^\vee"] \& \& S^{\vee\leftrightharpoons}
\end{tikzcd}}

Ceci posé, \deleted{\unsure{on a} $\unsure{\longrightarrow}$} la donnée d'une \deleted{sous $Z$-torseur $\pi$ de $\Pi_Z$}
\inserted{$Z$-jauge stricte} peut \unsure{s'interpréter|s'interprêter} en termes
de
\[
\left\{\begin{aligned}
 &\text{a) une }\deleted{\unsure{\text{iso}}}\text{ bijection }S^\leftrightharpoons\xrightarrow{\varphi^\leftrightharpoons}S^{\vee\leftrightharpoons}
 \quad\text{qui soit}\\
 &\hspace{13em}\text{une }R^\leftrightharpoons\text{-jauge}\\
 &\deleted{\text{b) un }Z\text{-torseur }\pi}\\
 &\text{b) un }\deleted{\text{torseur}}\inserted{\text{isomorphisme}}\\
 &\deleted{\unsure{g}}(S/S^\leftrightharpoons)\wedge_{Z_{S^\leftrightharpoons}}(\varphi^{\leftrightharpoons *}(S^\vee/S^{\vee\leftrightharpoons}))
 \simeq 1_{\deleted{\unsure{Z}}S^\leftrightharpoons}
\end{aligned}\right.\tag*{(48.6)}
\]
\margin{(i.e. donnant lieu à une relation \uline{symétrique} sur
$S^\leftrightharpoons\times S^\leftrightharpoons$, \deleted{\unsure{soit} $Z\unsure{R}$} soit
$\unsure{R^\leftrightharpoons_{S^\leftrightharpoons}}$}
% slanted
\margin{\deleted{\unsure{la précisant} : la donnée d'une \unsure{bijection} qui soit \unsure{anticommutant} à l'opération de $Z$, i.e satisfasse $\varphi\circ\lambda=\lambda^{-1}\circ\varphi$}}
\deleted{ou autrement dit la donnée d'une \unsure{de}}
% archive sheet number: 233

